
% Authoritative LaTeX manuscript; Markdown is a generated compatibility projection.

\emph{What justifies this claim?}

Start with the claim and identify the evidence capable of supporting it. Different properties require different evidence: a structural invariant may be checked exhaustively, a performance claim may require measurement, a semantic-fidelity claim may require generative evaluation and human judgment, and a formal proof establishes only the property encoded under its stated model and assumptions.

For each claim, ask four questions:

\begin{itemize}
\tightlist
\item
  \textbf{Claim.} What exactly is being asserted, and over what scope?
\item
  \textbf{Evidence.} What observation, measurement, proof, test, or other artifact bears on that claim?
\item
  \textbf{Evaluation.} What procedure interprets the evidence, and what can that procedure miss?
\item
  \textbf{Consequence.} What consequence should the result have? Evidence may inform a human, produce a warning, or justify a blocking gate; evidentiary strength and enforcement remain separate decisions.
\end{itemize}

% mage-code-language: text
\begin{verbatim}
  EVIDENCE QUALITY                            "What justifies this claim?"
  ──────────────────────────────────────────────────────────────────────
  CLAIM
    ↓  what would settle it?
  EVIDENCE
    ↓  how is it interpreted?
  EVALUATION
    ↓  what does the result deserve?
  CONSEQUENCE
    advisory · warning · validator result · gate · human decision
  ──────────────────────────────────────────────────────────────────────
  No backing      →  unsupported claim
  Weak backing    →  state the limit
  Strong backing  →  still choose the consequence deliberately
\end{verbatim}

The critical failure is claiming more confidence than the evidence supports when the claim matters. Preserve evidence provenance, state limitations explicitly, and distinguish what the evidence establishes from the consequence the organization chooses to attach.
